% Example output — Dr. Jordan Chen cover letter for Whitfield University
% Generated by claude-resume-kit for demonstration purposes
% This is a fictional researcher; all data is fabricated.

\documentclass[11pt,a4paper,roman]{moderncv}
\usepackage[english]{babel}
\moderncvstyle{classic}
\moderncvcolor{green}
\usepackage[utf8]{inputenc}
\usepackage{ragged2e}
\usepackage[scale=0.79]{geometry}
\usepackage[version=4,arrows=pgf-filled]{mhchem}
\renewcommand*{\makeletterclosing}{\par\vspace{2ex}\closingname\par}

\name{Jordan}{Chen}
\address{Richland, WA 99354}
\phone[mobile]{+1 (555) 123-4567}
\email{jordan.chen@email.com}

\begin{document}

\recipient{To}{Hiring Committee\\Department of Biomedical Engineering\\Whitfield University\\Westbrook, MA 02481}
\date{\today}
\opening{Dear Members of the Hiring Committee,}
\makelettertitle

\begin{justify}
Your department's work at the intersection of structural biology and therapeutic design resonates with the research program I have built over the past three years: using machine learning to accelerate protein engineering decisions that would otherwise take months of experimental iteration. As a postdoctoral researcher developing ML-guided enzyme screening pipelines, I am excited to apply for the Assistant Professor position in Computational Protein Engineering (BME-2026-0042), where I would establish an independent group bridging data-driven protein design and biomolecular simulation.

At Lakewood University, I fine-tuned the ESM-2 protein language model on 45,000 experimental melting temperatures to screen 8,500 enzyme variants for industrial thermostability, compressing what would have been over a year of wet-lab work into 48 hours of computation. Five of our top seven candidates were confirmed experimentally by collaborators via differential scanning calorimetry. I then co-developed an open-source transfer learning framework that reduces labeled training data requirements by 60\%, now adopted by four external research groups. More recently, I extended our ML pipeline to predict enzyme tolerance across eight organic co-solvent systems, opening a pathway toward engineering biocatalysts for green chemistry. This research trajectory, from classical MD to ML-accelerated protein engineering, reflects the kind of program I would build at Whitfield.

My doctoral work at Westfield established the simulation foundations that make this ML approach rigorous rather than purely correlative. I developed enhanced sampling protocols that predict protein folding temperatures within 8 K of experiment, benchmarked four force fields for intrinsically disordered proteins, and calculated ligand binding free energies across three drug target families with sub-kcal/mol accuracy. I also built the curated thermostability database that directly enabled my postdoctoral ML work. Throughout, I mentored three graduate students and developed computational lab modules now used department-wide.

I would welcome the opportunity to contribute to Whitfield's strengths in biomaterials and therapeutic design. I could teach courses in computational biology, molecular modeling, and machine learning for biomedical applications. I look forward to discussing how my research program would complement your department's existing strengths and the collaborative opportunities available through your HPC infrastructure.
\end{justify}

\vspace{0.3cm}
{Sincerely,\\
Jordan Chen, Ph.D.\\
Postdoctoral Research Associate\\
Lakewood University}

\end{document}
